\documentclass{article}
\usepackage{amsmath,amssymb}
\usepackage{xurl}
\usepackage[hidelinks]{hyperref}
% Shared commands for notes in docs/paper-gaps/.

\DeclareUnicodeCharacter{2080}{\ifmmode{}_{0}\else\textsubscript{0}\fi}
\DeclareUnicodeCharacter{2081}{\ifmmode{}_{1}\else\textsubscript{1}\fi}
\DeclareUnicodeCharacter{2082}{\ifmmode{}_{2}\else\textsubscript{2}\fi}
\DeclareUnicodeCharacter{2099}{\ifmmode{}_{n}\else\textsubscript{n}\fi}

\newcommand{\Lean}{\textsc{Lean}}
\ExplSyntaxOn
\NewDocumentCommand{\leanid}{m}{
  \ifmmode
    \textsf{\detokenize{#1}}
  \else
    \group_begin:
    \urlstyle{sf}
    \tl_set:Nn \l_tmpa_tl {#1}
    \tl_replace_all:Nnn \l_tmpa_tl {\_} {_}
    \exp_args:NV \path \l_tmpa_tl
    \group_end:
  \fi
}
\ExplSyntaxOff
\providecommand{\tr}{\operatorname{tr}}
\newcommand{\tnleanIssueBase}{https://github.com/LionSR/TNLean/issues/}
\newcommand{\tnleanPrBase}{https://github.com/LionSR/TNLean/pull/}
\newcommand{\tnleanDocBase}{https://sirui-lu.com/TNLean/docs/}
\newcommand{\ghissue}[1]{\href{\tnleanIssueBase#1}{GitHub issue \##1}}
\newcommand{\ghpr}[1]{\href{\tnleanPrBase#1}{PR \##1}}
\newcommand{\doclink}[1]{\href{\tnleanDocBase#1.html}{\path{#1}}}

% Dirac notation and the identity operator, as in blueprint/src/macros/common.tex.
% \providecommand keeps note-local definitions authoritative.
\usepackage{dsfont}
\providecommand{\ket}[1]{|#1\rangle}
\providecommand{\bra}[1]{\langle #1|}
\providecommand{\braket}[2]{\langle #1 | #2 \rangle}
\providecommand{\ketbra}[2]{|#1\rangle\!\langle #2|}
\providecommand{\Id}{\mathds{1}}
\providecommand{\one}{\mathds{1}}
\providecommand{\C}{\mathbb{C}}
\providecommand{\MN}[1]{M_{#1}(\C)}
\providecommand{\MD}{M_D(\C)}

% Verdict marker: \gapnote{<kind>}{<status>}. Kind is the policy
% classification (clarification, local-correction, scope-restriction,
% unfaithful, false-source, open-gap); status is open, wip (elimination
% underway), resolved, or historical. Machine-read by the site index;
% prints a verdict line.
\newcommand{\gapnote}[2]{\par\noindent\textbf{Verdict:} #1 (#2).\par}

\title{The Constant in the Depth Lower Bound}
\author{}
\date{2026-10-03}
\begin{document}
\maketitle
\gapnote{local-correction}{resolved}

\section{What the source states}
Let $\ket{\phi_N}$ be the normalized translation-invariant matrix product
states on $N$ sites generated by a normal tensor $A$ with correlation length
$\xi=-1/\ln\lvert\lambda_2\rvert>0$, and let $\ket{\psi_N}$ be obtained from
product states by local circuits of depth $T$. Theorem~1
of~\cite{Malz2024Preparation} states that if $T=o(\log N)$, then
$\varepsilon(\phi_N,\psi_N)=1-\lvert\braket{\phi_N}{\psi_N}\rvert>1/2$ for all
large $N$. Its proof is quantitative.\footnote{Source traceability: section
    ``Proof of Theorem~1'' of the Supplemental Material, the last two
    paragraphs of the proof; source at
    \url{https://arxiv.org/abs/2307.01696}.} With
$\mu=c_1e^{-2T/\xi}$, $\delta>\mu^2/9$ and $k=\lfloor N/(4q)\rfloor$ blocks of
length $q=\lceil2(1+\eta)\xi\ln N\rceil$, it asserts that $k>2/\delta$ gives
$(1-\delta)^{k/2}<1/e$ ``and thus $d(\phi,\psi)>\sqrt{3/4}$'', and closes with
the chain
\begin{align}
    k>\frac{N}{5q}>\frac{N\gamma}{10\xi\log N}>\frac{18}{c_1^2}e^{4T/\xi}
    >\frac2\delta .
    \label{eq:mswc24_depth_lower_bound_constant_chain}
\end{align}

\section{What fails}
Two constants in this passage are wrong; neither affects the rate $4/\xi$.
\begin{enumerate}
    \item From $(1-\delta)^{k/2}<1/e$ the bound
          $d(\phi,\psi)>1-(1-\delta)^{k/2}-\sqrt{2c_0}N^{-\eta}$ gives
          $d(\phi,\psi)>1-1/e-o(1)\approx0.632$, not $d>\sqrt{3/4}\approx0.866$.
          The conclusion $\varepsilon>1/2$ needs $\sqrt{3/4}$, because
          $d(\phi,\psi)=\sqrt{1-\lvert\braket{\phi}{\psi}\rvert^2}$ for pure
          states. It is reached once $(1-\delta)^{k/2}<1-\sqrt{3/4}-o(1)$, for
          example once $k>5/\delta$, since
          $2\ln\bigl(1/(1-\sqrt{3/4})\bigr)<4.1$.
    \item The factor $\gamma$ in
          (\ref{eq:mswc24_depth_lower_bound_constant_chain}) is not defined in
          the proof. With $q=\lceil2(1+\eta)\xi\ln N\rceil$ one has
          $N/(5q)\ge N/(10(1+\eta)\xi\ln N)$ up to the rounding of $q$, so
          $\gamma$ stands for a constant close to $1/(1+\eta)$.
\end{enumerate}
The factor $\tfrac12$ missing from the trace-distance bound, recorded
in
\path{docs/paper-gaps/mswc24_decaying_correlations_windowed_connected.tex}, also changes
only $c_1$. After these repairs the chain still reads
$N/\ln N\ge K\,e^{4T/\xi}$ with $K$ depending only on $A$; so
$\varepsilon(\phi_N,\psi_N)\le1/2$ forces
\begin{align}
    T\ge\frac\xi4\bigl(\ln N-\ln\ln N\bigr)-C_A .
    \label{eq:mswc24_depth_lower_bound_constant_source}
\end{align}
The source's argument does not give $T\ge(\xi/4)\ln N-C_A$ without the
correction $\ln\ln N$, which comes from the block length $q\propto\ln N$.

\section{What is formalized}
The formal bound is derived from the averaging argument of the blueprint,
which replaces the source's fidelity of blocked product states (see
\path{docs/paper-gaps/mswc24_depth_lower_bound_gauge.tex}), and carries the same rate. For a
normal tensor in the gauge of the source's eq.~(5), a unit vector $\psi$
prepared in depth $T$ with $\lvert\braket{\phi_N}{\psi}\rvert\ge1/2$ and
$N\ge B(T+1)$ satisfies
\begin{align}
    N\le C\,(T+1)\,e^{4T/\xi},
    \label{eq:mswc24_depth_lower_bound_constant_formal}
\end{align}
with $B,C$ independent of $N$ and $T$. The exponent is the square of the
correlator bound $c\,e^{-(s'-1)/\xi}$ at a separation $s'\le2T+O(1)$, and the
factor $T+1$ comes from the spacing of the windows, as the factor $\ln N$ in
(\ref{eq:mswc24_depth_lower_bound_constant_chain}) comes from the block
length. Taking logarithms, $\varepsilon(\phi_N,\psi)\le1/2$ forces
$(\xi/4)\ln N\le T+(\xi/4)\ln(T+1)+C'$ for every normal tensor, with $\xi$
computed from $\lambda_2/\lambda_1$; since $T<(\xi/4)\ln N$ otherwise, this
implies (\ref{eq:mswc24_depth_lower_bound_constant_source}). For every
$a<\xi/4$ it gives $T\ge a\ln N-C_a$, and with the source's eq.~(1) the
optimal depth at fixed error $0<\varepsilon\le1/2$ is
$\Theta(\log N)$.\footnote{Formal declarations:
    \leanid{MPSTensor.exists_depth_lower_bound} in
    \doclink{TNLean/MPS/Preparation/DepthLowerBound};
    \leanid{MPSTensor.exists_log_le_of_infidelity_le_one_half},
    \leanid{MPSTensor.exists_log_le_of_infidelity_le_one_half_of_isNormal},
    \leanid{MPSTensor.exists_mul_log_le_of_infidelity_le_one_half_of_isNormal}
    and \leanid{MPSTensor.exists_depth_isTheta_log} in
    \doclink{TNLean/MPS/Preparation/DepthLowerBoundConstant}.}

\section{Verdict}
The two slips are in constants of the closing inequality; the theorem and
the rate $4/\xi$, that is the constant $\xi/4$ in front of $\ln N$, stand.
The formal statements carry the source's rate with the logarithmic
correction written as $(\xi/4)\ln(T+1)$.

\bibliographystyle{alpha}
\bibliography{references}
\end{document}
